% =====================================================================
%  Thème 4 — Produit scalaire et orthogonalité — MOYEN — Exercices
% =====================================================================
\documentclass[11pt,a4paper]{article}
\usepackage[utf8]{inputenc}
\usepackage[T1]{fontenc}
\usepackage{lmodern}
\usepackage[french]{babel}
\usepackage{amsmath,amssymb,mathtools}
\usepackage{xcolor}
\usepackage{tikz}
\usepackage{enumitem}
\usepackage{fancyhdr}
\usepackage[a4paper,margin=2.1cm,headheight=26pt,headsep=10pt]{geometry}
\usetikzlibrary{arrows.meta,calc}

\definecolor{primary}{RGB}{222,70,80}
\definecolor{primarydark}{RGB}{150,30,42}
\definecolor{axiscol}{RGB}{90,90,90}
\definecolor{vecu}{RGB}{20,20,20}

\pagestyle{fancy}\fancyhf{}
\fancyhead[L]{\small\textcolor{primarydark}{\textbf{Fiche 2} — Calcul vectoriel}}
\fancyhead[R]{\small\textcolor{primarydark}{\textsc{Thème 4 — Moyen — Exercices}}}
\fancyfoot[C]{\small\thepage}
\renewcommand{\headrulewidth}{0.8pt}

\newcommand{\vc}[1]{\overrightarrow{#1}}
\providecommand{\norm}[1]{\left\lVert #1\right\rVert}
\newcommand{\exo}[1]{\medskip\noindent{\color{primary}\bfseries\large Exercice #1.}\par\nopagebreak\smallskip}

\begin{document}
\begin{center}
{\LARGE\bfseries\color{primarydark}Produit scalaire et orthogonalité — Niveau Moyen}\\[2pt]
{\color{primary}\rule{0.55\linewidth}{1.1pt}}
\end{center}
\vspace{1mm}
{\small\itshape Objectif : imposer l'orthogonalité pour trouver une inconnue, calculer un produit
scalaire à partir de points, utiliser la linéarité et exploiter une figure.}
\vspace{2mm}

\exo{1} \textbf{Trouver une composante par orthogonalité.}
Déterminer la (ou les) valeur(s) rendant $\vec u$ et $\vec v$ orthogonaux.
\begin{enumerate}[label=\textbf{\alph*)},leftmargin=2em,itemsep=2pt]
  \item $\vec u\,(3;y)$ et $\vec v\,(2;6)$.
  \item $\vec u\,(x;4)$ et $\vec v\,(2;-1)$.
  \item $\vec u\,(m;2)$ et $\vec v\,(m;-8)$.
\end{enumerate}

\exo{2} \textbf{Produit scalaire à partir de points.}
On donne $A\,(1;1)$, $B\,(4;2)$ et $C\,(2;5)$.
\begin{enumerate}[label=\textbf{\alph*)},leftmargin=2em,itemsep=2pt]
  \item Calculer les composantes de $\vc{AB}$ et $\vc{AC}$.
  \item Calculer $\vc{AB}\cdot\vc{AC}$. L'angle $\widehat{BAC}$ est-il aigu, droit ou obtus ?
\end{enumerate}

\exo{3} \textbf{Droites perpendiculaires par leurs équations.}
On considère $d:\ 2x-3y+1=0$ et $d':\ 3x+2y-4=0$.
\begin{enumerate}[label=\textbf{\alph*)},leftmargin=2em,itemsep=2pt]
  \item Donner un vecteur directeur de $d$ et un vecteur directeur de $d'$.
  \item Montrer que $d$ et $d'$ sont perpendiculaires.
\end{enumerate}

\exo{4} \textbf{Utiliser la linéarité.}
On donne $\vec u\,(2;3)$, $\vec v\,(1;-1)$ et $\vec w\,(4;0)$.
\begin{enumerate}[label=\textbf{\alph*)},leftmargin=2em,itemsep=2pt]
  \item Calculer $\vec u\cdot\vec v$ et $\vec w\cdot\vec v$.
  \item En déduire $(2\vec u)\cdot\vec v$ puis $(\vec u+\vec w)\cdot\vec v$, et vérifier par un calcul direct.
\end{enumerate}

\exo{5} \textbf{Produit scalaire dans un carré.}
$ABCD$ est le carré de sommets $A\,(0;0)$, $B\,(2;0)$, $C\,(2;2)$, $D\,(0;2)$.
\begin{enumerate}[label=\textbf{\alph*)},leftmargin=2em,itemsep=2pt]
  \item Calculer $\vc{AB}\cdot\vc{AD}$. Que traduit ce résultat ?
  \item Calculer $\vc{AB}\cdot\vc{AC}$.
  \item Calculer $\vc{AC}\cdot\vc{BD}$ et interpréter géométriquement (diagonales du carré).
\end{enumerate}

\end{document}
